% Бүлэг 1

\chapter{Төслийн түүх, холбоотой судалгаа ба систем хөгжүүлэх арга зүй}
\label{Chapter1}

\section{Бүлгийн зорилго ба судалгааны хүрээ}
Энэхүү бүлэг нь дипломын ажлын эхний судалгааны суурь хэсэг бөгөөд системийн бодит хэрэгжилт рүү шууд орохоос өмнө төслийн үүсэл, хөгжлийн шалтгаан, ижил төстэй шийдлүүд, хэрэглэгчийн хэрэгцээ, программ хангамж хөгжүүлэх арга зүйн сонголтыг тайлбарлана. Өөрөөр хэлбэл, энэ бүлгийн үндсэн зорилго нь \textit{“яагаад ийм систем хэрэгтэй вэ, ямар төрлийн системүүдтэй хамааралтай вэ, ямар хөгжүүлэлтийн аргачлалаар хэрэгжүүлэх нь зохистой вэ?”} гэсэн гурван асуултад академик байдлаар хариулахад оршино.

Энэхүү төсөл нь FL Studio болон хөгжим боловсруулалт суралцагчдад зориулсан веб сургалтын платформ юм. Систем нь курсийн каталог, төлбөртэй хичээлийн хандалт, хичээлийн тоглуулагч, өөрийгөө шалгах сорил, суралцах явцын хяналт, админ удирдлага, мөн RAG суурьтай AI туслахыг нэгтгэнэ. Гэхдээ эдгээр модулийн техникийн хэрэгжилтийг энэ бүлэгт дэлгэрэнгүй авч үзэхгүй. Харин тэдгээрийг хөгжүүлэхэд ямар судалгааны үндэслэл, ямар арга зүйн сонголт хэрэгтэй байсныг тодорхойлно.

Судалгааны хүрээг дараах байдлаар тогтоов.
\begin{itemize}
  \item төслийн хөгжлийн үе шат, гол шийдвэрүүдийг нэгтгэн тайлбарлах;
  \item онлайн сургалтын платформ, хөгжмийн сургалтын систем, AI туслахтай сургалтын орчны ижил төстэй шийдлүүдийг харьцуулах;
  \item хэрэглэгчийн хэрэгцээ, функциональ болон функциональ бус шаардлагын эхний түвшний шинжилгээ хийх;
  \item Waterfall, Prototype, Spiral, Agile/Scrum, Kanban, Lean UX, Design Science Research зэрэг аргачлалуудыг харьцуулж, төслийн нөхцөлд тохирох арга зүйг сонгох;
  \item сонгосон арга зүйг энэ дипломын төслийн бодит ажлын төлөвлөгөөтэй уялдуулах.
\end{itemize}

\section{Төслийн үүсэл, хөгжлийн түүх}
Энэхүү дипломын төсөл нь анх хиймэл оюун ухааны тусламжтай хөгжим боловсруулалтын вебсайт бүтээх санаанаас эхэлсэн. Хөгжүүлэлтийн явцад уг санаа зөвхөн нэг AI функцтэй туршилтын систем бус, харин сургалтын контент, төлбөртэй хандалт, суралцах явц, өгөгдлийн сан, AI туслахыг нэгтгэсэн сургалтын платформ болж өргөжсөн. Репозиторийн бүтцэд \texttt{Frontend}, \texttt{dataset}, \texttt{rag-dataset-builder}, \texttt{diplomapaper} зэрэг бүрэлдэхүүнүүд зэрэгцэн оршиж байгаа нь төслийг нэг бүтээгдэхүүн, нэг судалгааны объект гэж авч үзэх үндэслэл болно.

\begin{figure}[htbp]
  \centering
  \resizebox{\textwidth}{!}{\begin{tikzpicture}[x=1cm,y=1cm,>=latex]
  \tikzstyle{phase}=[draw, rounded corners=2pt, fill=blue!8, minimum width=3.2cm, minimum height=1.1cm, align=center]
  \tikzstyle{arrow}=[->, thick]

  \node[phase] (p1) at (0,0) {1-р үе шат\\Контент ба туршлагын эхлэл};
  \node[phase] (p2) at (4.2,0) {2-р үе шат\\Курс, хичээлийн бүтэц};
  \node[phase] (p3) at (8.4,0) {3-р үе шат\\Үүлэн сангийн интеграц};
  \node[phase] (p4) at (12.6,0) {4-р үе шат\\Мэдээлэл сэргээх чатбот ба админ};

  \draw[arrow] (p1) -- (p2);
  \draw[arrow] (p2) -- (p3);
  \draw[arrow] (p3) -- (p4);

  \node[align=center] at (0,-1.2) {Нүүр хуудас\\Статик өгөгдөл};
  \node[align=center] at (4.2,-1.2) {Хөтөлбөр\\Өөрийгөө шалгах};
  \node[align=center] at (8.4,-1.2) {Нэвтрэлт, худалдан авалт\\ахиц};
  \node[align=center] at (12.6,-1.2) {Чатны зам\\Админ самбар};
\end{tikzpicture}
}
  \caption[Төслийн хөгжлийн үе шат]{Төслийн хөгжлийн үндсэн үе шат ба бүрэлдэхүүнүүдийн уялдаа}
  \label{fig:project-timeline}
\end{figure}

\begin{longtable}{|p{0.13\textwidth}|p{0.22\textwidth}|p{0.30\textwidth}|p{0.25\textwidth}|}
  \caption{Төслийн хөгжлийн үе шат ба судалгааны утга}
  \label{tab:project-history}\\
  \hline
  \textbf{Үе шат} & \textbf{Гол зорилго} & \textbf{Бодит материал} & \textbf{Судалгааны ач холбогдол} \\ \hline
  I үе шат & Сургалтын бүтээгдэхүүний анхны чиглэлийг тодорхойлох & Landing page, hero болон learning path компонентууд & Суралцагчид зориулсан үнэ цэнийн санал, анхны UX урсгал, бүтээгдэхүүний байршуулалт тодорсон. \\ \hline
  II үе шат & Курс, хичээл, дадлагын бүтцийг бүрдүүлэх & Course data, catalog page, course detail page & Сургалтын контентыг курс, lesson, exercise, takeaway хэлбэрээр шаталсан загварт оруулсан. \\ \hline
  III үе шат & Хэрэглэгч, төлбөр, явцын өгөгдлийн суурь тавих & Supabase schema, auth болон purchase hooks & Бүртгэл, худалдан авалт, явц хадгалалт зэрэг платформын үндсэн бизнес логик тодорхой болсон. \\ \hline
  IV үе шат & AI туслах болон мэдлэгийн санг системд холбох & Chat API, RAG builder, dataset & RAG суурьтай асуулт-хариултын туслахыг сургалтын орчинд нэмэлт функц байдлаар нэгтгэсэн. \\ \hline
  V үе шат & Дипломын баримт бичгийн зохион байгуулалт & \texttt{diplomapaper} хавтас, диаграм, бүлгүүд & Бүтээгдэхүүний хэрэгжилтийг судалгаа, арга зүй, зохиомж, архитектуртай холбон тайлбарлах шаардлага үүссэн. \\ \hline
\end{longtable}

Энэ хөгжлийн түүхээс харахад төсөл нь эхнээсээ бүрэн тодорхой шаардлагатай, нэг удаагийн шугаман хөгжүүлэлтээр дуусах систем биш байв. Сургалтын контентын бүтэц, хэрэглэгчийн урсгал, төлбөртэй хандалт, AI туслахын үүрэг зэрэг нь хөгжүүлэлтийн явцад туршилт, засвар, дахин зохион байгуулалтаар боловсорсон. Иймээс арга зүйн хувьд давталттай, уян хатан, эрсдэлийг үе шаттай бууруулдаг хандлага шаардлагатай болсон.

\section{Судалгааны асуудал}
Хөгжим боловсруулалт сурах үйл явц нь програм ашиглах техникийн мэдлэг, хөгжмийн онолын ойлголт, бүтээлч дадлага, сонсголын үнэлгээ, микс ба мастерингийн шийдвэрийг зэрэг шаарддаг. Эхлэн суралцагчид ихэвчлэн YouTube, форум, богино нийтлэл, курсийн хэсэгчилсэн материал зэрэг олон эх сурвалжаас суралцдаг боловч тэдгээр эх сурвалжууд хоорондоо дараалал, түвшин, нэр томьёо, дадлагын аргачлалаараа зөрүүтэй байдаг.

Энэхүү асуудлыг гурван түвшинд авч үзэж болно.
\begin{enumerate}
  \item \textbf{Сургалтын зохион байгуулалтын асуудал:} эхлэн суралцагч ямар хичээлээс эхлэх, дараагийн алхамд ямар чадвар эзэмших, ямар дадлага хийхээ тодорхой мэдэхгүй.
  \item \textbf{Дадлага ба ахицын хэмжилтийн асуудал:} видео үзсэн эсэхээс гадна тухайн ойлголтыг хэрхэн шалгах, ямар дасгалаар бататгах, ахицаа хэрхэн хадгалах нь тодорхойгүй.
  \item \textbf{Тухайн мөчийн зөвлөгөөний асуудал:} суралцагч хичээл үзэж байхдаа ойлгоогүй нэр томьёо, plugin тохиргоо, mixing шийдвэр, FL Studio-ийн үйлдэлтэй холбоотой жижиг асуултыг тэр дор нь тодруулах хэрэгцээтэй.
\end{enumerate}

Эдгээр асуудлыг зөвхөн AI чатботоор шийдвэрлэх боломжгүй. AI нь буруу, ерөнхий эсвэл эх сурвалжгүй зөвлөгөө өгөх эрсдэлтэй тул үндсэн сургалтын контентыг орлох бус, харин бүтэцтэй сургалтын урсгалыг дэмжих хэлбэрээр ашиглагдах ёстой. Иймээс төслийн үндсэн судалгааны асуудлыг дараах байдлаар тодорхойлов.

\begin{quote}
FL Studio болон хөгжим боловсруулалт суралцагчдад зориулсан сургалтын платформд курсийн бүтэц, төлбөртэй хандалт, ахицын хяналт, AI туслахыг хэрхэн уялдаа холбоотой, өргөтгөх боломжтой, хэрэглэгч төвтэй байдлаар хөгжүүлэх вэ?
\end{quote}

\section{Хэрэглэгч ба оролцогч талуудын шинжилгээ}
Системийн үндсэн хэрэглэгч нь FL Studio болон хөгжим боловсруулалт шинээр сурах, эсвэл эхний түвшний чадвараа цэгцлэх хүсэлтэй суралцагч юм. Үүнээс гадна багш, продюсер, админ, систем хөгжүүлэгч, контент бэлтгэгч зэрэг оролцогч талууд системд шууд болон шууд бусаар нөлөөлнө.

\begin{table}[htbp]
  \centering
  \caption{Оролцогч талууд ба хэрэгцээ}
  \label{tab:stakeholders}
  \begin{tabular}{|p{0.18\textwidth}|p{0.34\textwidth}|p{0.36\textwidth}|}
    \hline
    \textbf{Оролцогч} & \textbf{Гол хэрэгцээ} & \textbf{Системийн хариу шийдэл} \\ \hline
    Зочин хэрэглэгч & Курсийн мэдээлэл үзэх, үнэ цэнэ ойлгох, эхний хичээлийг турших & Landing page, course catalog, free lesson preview \\ \hline
    Суралцагч & Курс үзэх, төлбөртэй хичээлд хандах, ахицаа хадгалах, асуулт асуух & Lesson player, purchase access, self-check quiz, progress tracking, AI mentor \\ \hline
    Багш/продюсер & Контентоо бүтэцтэй хүргэх, төлбөртэй бүтээгдэхүүн болгох & Course, lesson, teacher profile, админ удирдлагын боломж \\ \hline
    Админ & Курс, хичээл, хэрэглэгч, мэдлэгийн санг хянах & Admin dashboard, Supabase өгөгдлийн загвар, RLS бодлого \\ \hline
    Систем хөгжүүлэгч & Модульчлагдсан, засварлахад ойлгомжтой бүтэцтэй кодын бааз & Next.js App Router, компонент бүтэц, API route, Supabase integration \\ \hline
  \end{tabular}
\end{table}

Хэрэглэгчийн урсгалын хувьд систем нь \textit{танилцах} → \textit{курс сонгох} → \textit{хандалт баталгаажуулах} → \textit{хичээл үзэх} → \textit{өөрийгөө шалгах} → \textit{AI туслахаас лавлах} → \textit{ахицаа хадгалах} гэсэн шатлалтай. Энэ нь цэвэр мэдээллийн сайт бус, харин суралцах үйл явцыг удирдах систем болохыг харуулна.

\section{Функциональ шаардлагын эхний шинжилгээ}
Функциональ шаардлага нь хэрэглэгч системээр ямар үйлдэл хийх ёстойг тодорхойлно. Энэ бүлэгт шаардлагыг ерөнхий түвшинд авч үзэх ба нарийвчилсан зохиомж, өгөгдлийн бүтэц, API урсгалыг дараагийн бүлгүүдэд тайлбарлана.

\begin{longtable}{|p{0.08\textwidth}|p{0.27\textwidth}|p{0.52\textwidth}|}
  \caption{Системийн үндсэн функциональ шаардлага}
  \label{tab:functional-requirements}\\
  \hline
  \textbf{Код} & \textbf{Шаардлага} & \textbf{Тайлбар} \\ \hline
  FR-01 & Хэрэглэгч бүртгүүлэх ба нэвтрэх & Суралцагч өөрийн бүртгэлээр системд нэвтэрч худалдан авалт, ахиц, чат түүх зэрэг хувийн мэдээлэлтэй ажиллана. \\ \hline
  FR-02 & Курс хайх, шүүх, эрэмбэлэх & Хэрэглэгч курсийг ангилал, түвшин, үнэ, хайлтын үгээр шүүж өөрт тохирох сургалтын замналаа сонгоно. \\ \hline
  FR-03 & Хичээлийн дэлгэрэнгүй үзэх & Курсийн тайлбар, багш, үнэ, free/paid lesson, curriculum, duration зэрэг мэдээллийг харуулна. \\ \hline
  FR-04 & Төлбөртэй хичээлийн хандалт шалгах & Систем хэрэглэгч тухайн курсийг худалдан авсан эсэхээс хамаарч lesson access-ийг нээж эсвэл хязгаарлана. \\ \hline
  FR-05 & Хичээл тоглуулах ба lesson brief үзүүлэх & Видео бүхий lesson-г тоглуулах, видео байхгүй үед summary, exercise, takeaway хэлбэрийн богино сургалтын материалыг харуулна. \\ \hline
  FR-06 & Өөрийгөө шалгах сорил & Хэрэглэгч lesson болон course түвшний асуултад хариулж, ойлголтоо бататгана. \\ \hline
  FR-07 & Суралцах явц хадгалах & Дуусгасан lesson, course progress, хамгийн сүүлд үзсэн lesson зэрэг мэдээллийг хадгална. \\ \hline
  FR-08 & AI туслахаас асуулт асуух & Суралцагч FL Studio, mixing, mastering, хөгжмийн онолын асуултаа чатботод илгээж RAG контексттэй хариу авна. \\ \hline
  FR-09 & Админ удирдлага & Админ курс, lesson, хэрэглэгчийн өгөгдөл, мэдлэгийн сан зэрэг мэдээллийг хянах, засах боломжтой байна. \\ \hline
  FR-10 & Dataset бэлтгэх ба импортлох & RAG мэдлэгийн сангийн өгөгдлийг JSONL/CSV хэлбэрээр бэлтгэж Supabase knowledge base-д импортлох боломжтой байна. \\ \hline
\end{longtable}

\section{Функциональ бус шаардлагын эхний шинжилгээ}
Функциональ бус шаардлага нь системийн чанарын шинжүүдийг тодорхойлно. Сургалтын платформын хувьд хэрэглэгчийн итгэл, өгөгдлийн хамгаалалт, хариу өгөх хурд, өргөтгөх боломж, засварлахад хялбар байдал чухал.

\begin{table}[htbp]
  \centering
  \caption{Функциональ бус шаардлагын ангилал}
  \label{tab:nonfunctional-requirements}
  \begin{tabular}{|p{0.22\textwidth}|p{0.30\textwidth}|p{0.34\textwidth}|}
    \hline
    \textbf{Чанарын шинж} & \textbf{Шаардлагын утга} & \textbf{Төслийн хүрээнд хэрэгжүүлэх чиглэл} \\ \hline
    Ашиглахад хялбар байдал & Анхан шатны суралцагч интерфейсийг тайлбаргүй ойлгох & Course catalog, lesson player, dashboard, чат товчийг ойлгомжтой байрлуулах \\ \hline
    Найдвартай байдал & Хичээлийн хандалт, ахиц, төлбөртэй эрх алдагдахгүй байх & Supabase өгөгдлийн сан, unique constraint, RLS бодлого ашиглах \\ \hline
    Аюулгүй байдал & Хэрэглэгч зөвхөн өөрийн өгөгдөлд хандах & Authentication, authorization, row level security \\ \hline
    Өргөтгөх чадвар & Шинэ курс, багш, dataset, AI provider нэмэхэд суурь бүтэц эвдрэхгүй байх & Модульчлагдсан компонент, API route, provider abstraction \\ \hline
    Гүйцэтгэл & Курсийн жагсаалт, хичээлийн хуудас, чат хариу боломжийн хугацаанд ажиллах & Next.js rendering, local dataset cache, provider fallback \\ \hline
    Академик найдвартай байдал & AI хариулт эх сурвалжгүйгээр албан ёсны мэдлэг мэт харагдахгүй байх & RAG context, citation-safe flag, generated guidance-ийг ялгах \\ \hline
  \end{tabular}
\end{table}

\section{Ижил төстэй системүүдийн судалгаа}
Ижил төстэй системүүдийг судлахдаа зөвхөн хөгжмийн AI бүтээгдэхүүнүүдийг бус, онлайн сургалтын платформ, хөгжмийн сургалтын орчин, AI туслахтай мэдлэгийн систем, DAW хэрэглэгчийн сургалтын эх сурвалжуудыг хамруулсан. Учир нь энэхүү дипломын төсөл эдгээр дөрвөн чиглэлийн уулзвар дээр байрлаж байна.

\subsection{Онлайн сургалтын платформууд}
Coursera, Udemy, Skillshare зэрэг платформууд курсийн каталог, төлбөртэй контент, багшийн профайл, хэрэглэгчийн явцын хяналт, үнэлгээ зэрэг нийтлэг функцтэй. Эдгээр системийн давуу тал нь сургалтын бүтээгдэхүүнийг худалдаалах, хичээлийг багцлах, суралцагчийн хэрэглээг хэмжих туршлага юм. Гэвч ерөнхий платформууд нь FL Studio-ийн тодорхой ажлын урсгал, монгол-англи холимог нэр томьёо, тухайн хичээлийн мөчид асуулт асуух AI туслахыг заавал агуулдаггүй.

\subsection{Хөгжмийн сургалтын эх сурвалжууд}
YouTube, FL Studio-ийн албан ёсны гарын авлага, продюсеруудын богино tutorial, хөгжмийн форумууд нь хөгжмийн продакшн сурахад өргөн хэрэглэгддэг. Эдгээр нь хүртээмжтэй боловч эхлэн суралцагчид хэрэгтэй дараалсан замнал, ахицын хэмжилт, төлбөртэй course access, нэгтгэсэн dashboard-оор дутмаг. Иймээс энэ төслийн платформ нь нээлттэй эх сурвалжуудыг бүрэн орлох бус, харин тэдгээрээс сурахад тулгардаг зохион байгуулалтын асуудлыг багасгах зорилготой.

\subsection{AI хөгжмийн системүүд}
Jukebox, MusicLM, MusicGen зэрэг судалгааны ажлууд болон Magenta, Suno, AIVA, SOUNDRAW зэрэг бүтээгдэхүүнүүд нь AI-г хөгжим үүсгэх, audio modeling, бүтээлч санаа гаргах чиглэлд ашиглаж байна \cite{jukebox_dhariwal_2020,musiclm_2023,musicgen_2023,magenta_site,suno_ai_app,aiva_ai,soundraw_home}. Гэхдээ энэ дипломын ажилд AI-ийн үүрэг нь хөгжмийг хэрэглэгчийн өмнөөс бүрэн үүсгэх бус, харин хүний багшийн сургалтын контентыг дэмжих зөвлөгчийн үүрэг юм. Энэ ялгаа нь системийн эрсдэл, зохиогчийн эрх, сургалтын чанарын талаас чухал.

\begin{longtable}{|p{0.17\textwidth}|p{0.25\textwidth}|p{0.26\textwidth}|p{0.22\textwidth}|}
  \caption{Ижил төстэй системүүдийн харьцуулалт}
  \label{tab:related-systems}\\
  \hline
  \textbf{Системийн төрөл} & \textbf{Давуу тал} & \textbf{Хязгаарлалт} & \textbf{Энэ төсөлд авсан санаа} \\ \hline
  Онлайн сургалтын платформ & Курсийн бүтэц, төлбөр, багшийн профайл, ахицын хяналт сайн хөгжсөн & Ерөнхий сэдэвт төвлөрдөг, FL Studio-ийн workflow-д нарийн тохироогүй & Course catalog, paid access, dashboard, progress tracking \\ \hline
  YouTube/tutorial эх сурвалж & Хүртээмж өндөр, олон жишээтэй, бодит хэрэглээтэй & Дараалалгүй, түвшин холилдсон, ахиц хадгалахгүй & Free preview, lesson brief, curated learning path \\ \hline
  DAW гарын авлага & Албан ёсны, нарийвчилсан, нэр томьёо зөв & Эхлэн суралцагчид хэт техникийн, дадлага багатай & Knowledge base-д эх сурвалжийн metadata хадгалах \\ \hline
  AI music generator & Бүтээлч санаа хурдан гаргах, генератив боломж өндөр & Сургалтын замнал, хүний багшийн тайлбар, ахицын үнэлгээ хангалтгүй & AI-г контент бүтээгч бус, сургалтын туслах болгон ашиглах \\ \hline
  RAG chatbot & Асуултад контексттэй хариулах, мэдлэгийн сангаар хязгаарлах боломжтой & Dataset чанар, эх сурвалжийн найдвартай байдал, hallucination эрсдэлтэй & RAG dataset, citation-safe flag, fallback logic \\ \hline
\end{longtable}

\section{Холбоотой судалгааны ажлууд}
Энэхүү дипломын төслийн судалгааны суурь нь хоёр хэсэгтэй. Нэгдүгээрт, программ хангамж хөгжүүлэх арга зүй, шаардлагын шинжилгээ, хэрэглэгч төвтэй дизайн, дизайн шинжлэх ухааны судалгааны ажил; хоёрдугаарт, боловсролын системд AI ашиглах, RAG, аудио AI, хөгжмийн мэдээлэл боловсруулах судалгаанууд юм. AI онолын нарийвчилсан тайлбарыг хоёрдугаар бүлэгт авч үзэх тул энэ хэсэгт зөвхөн төслийн арга зүйтэй шууд холбоотой ажлуудыг нэгтгэнэ.

\begin{longtable}{|p{0.05\textwidth}|p{0.25\textwidth}|p{0.28\textwidth}|p{0.30\textwidth}|}
  \caption{Арга зүй ба систем хөгжүүлэлтийн холбоотой эх сурвалжууд}
  \label{tab:method-related-works}\\
  \hline
  \textbf{\#} & \textbf{Эх сурвалж} & \textbf{Гол санаа} & \textbf{Төслийн хэрэглээ} \\ \hline
  1 & Royce-ийн Waterfall загвар \cite{royce_waterfall_1970} & Том системд үе шат, баримтжуулалт, хяналтыг тодорхойлсон & Дипломын тайлан, шаардлага, зохиомжийн баримтжуулалтад хэрэгтэй боловч өөрчлөлт ихтэй веб системд дангаар тохиромжгүй. \\ \hline
  2 & Boehm-ийн Spiral загвар \cite{boehm_spiral_1988} & Эрсдэлд суурилсан давталттай хөгжүүлэлт & AI туслах, dataset чанар, provider сонголтын эрсдэлийг үе шаттай бууруулахад тохиромжтой. \\ \hline
  3 & Agile Manifesto \cite{agile_manifesto_2001} & Ажиллаж буй бүтээгдэхүүн, өөрчлөлтөд хариу үйлдэл, хэрэглэгчийн хамтын ажиллагааг чухалчилсан & Landing, course, dashboard, chat зэрэг модулийг богино давталтаар сайжруулах үндэс болсон. \\ \hline
  4 & Scrum Guide \cite{scrum_guide_2020} & Sprint, backlog, incremental delivery зарчим & Дипломын ажлын feature backlog, iteration, review төлөвлөлтөд ашиглах боломжтой. \\ \hline
  5 & Kanban аргачлал \cite{anderson_kanban_2010} & Work-in-progress хязгаарлах, урсгалыг харагдуулах & Нэг хүн эсвэл жижиг багийн дипломын төсөлд task flow-ийг хянахад тохиромжтой. \\ \hline
  6 & Lean UX \cite{lean_ux_2013} & Таамаглал, prototype, user feedback-д суурилсан UX хөгжүүлэлт & Сургалтын интерфейс, course discovery, chat drawer зэрэг хэрэглэгчийн урсгалыг туршихад хэрэгтэй. \\ \hline
  7 & Design Science Research \cite{hevner_dsr_2004} & Артефакт бүтээж, түүний хэрэгцээ ба үр нөлөөг үнэлэх судалгааны арга & Энэхүү дипломын ажил нь судалгаа + бодит платформ гэсэн хоёр үр дүнтэй тул DSR логиктой нийцнэ. \\ \hline
\end{longtable}

\section{Систем хөгжүүлэх аргачлалуудын харьцуулалт}
Программ хангамжийн аргачлал сонгохдоо төслийн хэмжээ, багийн бүтэц, шаардлагын тогтвортой байдал, эрсдэлийн төрөл, хэрэглэгчийн оролцоо, баримтжуулалтын хэрэгцээг харгалзан үздэг. Энэ дипломын төсөлд шаардлага хөгжүүлэлтийн явцад нарийссан, AI болон dataset хэсэг туршилт ихтэй, frontend UX олон давталтаар сайжирсан тул хатуу шугаман аргачлал дангаараа хангалтгүй.

\subsection{Waterfall аргачлал}
Waterfall аргачлал нь шаардлага, шинжилгээ, зохиомж, хэрэгжилт, туршилт, нэвтрүүлэлт гэсэн дараалсан үе шаттай \cite{royce_waterfall_1970}. Давуу тал нь баримтжуулалт сайн, хяналтын цэг тодорхой, сургалтын байгууллагын тайлангийн бүтэцтэй нийцдэг. Харин сул тал нь шаардлага өөрчлөгдөхөд засварын өртөг өсдөг. Энэхүү төслийн AI туслах, RAG dataset, хэрэглэгчийн урсгал зэрэг нь хөгжүүлэлтийн явцад тодорсон тул Waterfall-ийг зөвхөн баримтжуулалт, тайлангийн логикт ашиглах нь зохистой.

\subsection{Prototype аргачлал}
Prototype аргачлал нь хэрэглэгчийн ойлголт бүрэн тодорхойгүй үед хурдан загвар гаргаж, түүнийгээ туршин сайжруулахад тохиромжтой. Энэ төслийн landing page, course catalog, lesson detail, chat drawer зэрэг интерфейсүүд prototype хандлагаар боловсорсон. Давуу тал нь хэрэглэгчийн туршлагыг эрт шалгах боломжтой. Сул тал нь prototype-ийг бүтээгдэхүүний эцсийн код болгож үлдээвэл архитектурын өр нэмэгдэх эрсдэлтэй.

\subsection{Spiral аргачлал}
Spiral загвар нь эрсдэлийг давталт бүрд тодорхойлж, шийдэл туршин, дараагийн мөчлөгт шилждэг \cite{boehm_spiral_1988}. AI-тэй системд hallucination, dataset чанар, provider availability, latency, cost зэрэг эрсдэл өндөр байдаг тул энэ аргачлал оновчтой. Жишээлбэл, RAG chatbot эхний шатанд локал JSONL retrieval ашиглаж, дараагийн шатанд Supabase pgvector embedding рүү шилжих боломжтойгоор төлөвлөгдсөн.

\subsection{Agile ба Scrum}
Agile аргачлал нь шаардлага өөрчлөгдөхөд хурдан хариу үйлдэл үзүүлэх, ажиллаж буй бүтээгдэхүүнийг үе шаттай хүргэх, хэрэглэгчийн саналд тулгуурлах зарчимтай \cite{agile_manifesto_2001}. Scrum нь үүнийг sprint, backlog, review, retrospective гэсэн илүү тодорхой процесс болгон хэрэгжүүлдэг \cite{scrum_guide_2020}. Энэхүү дипломын төсөлд багийн хэмжээ бага боловч feature backlog, богино iteration, incremental delivery зарчим шууд хэрэглэгдэх боломжтой.

\subsection{Kanban}
Kanban нь ажлын урсгалыг харагдуулах, зэрэг эхлүүлсэн ажлыг хязгаарлах, тасралтгүй сайжруулахад төвлөрдөг \cite{anderson_kanban_2010}. Дипломын ажлын нөхцөлд frontend, database, dataset, thesis бичвэр зэрэг олон төрлийн ажил зэрэгцдэг тул Kanban самбар ашиглан \textit{Backlog → In progress → Review → Done} хэлбэрээр хянах нь тохиромжтой.

\subsection{Lean UX ба хэрэглэгч төвтэй хандлага}
Lean UX нь эхлээд том баримт бичиг биш, харин хэрэглэгчийн асуудлын таамаглал, хурдан prototype, туршилт, суралцсан зүйлд тулгуурладаг \cite{lean_ux_2013}. Хөгжмийн сургалтын платформын хувьд хэрэглэгч тухайн lesson-оос юу ойлгох, ямар үед AI туслахаас асуух, course card дээр ямар мэдээлэл харвал шийдвэр гаргах зэрэг нь UX туршилтаар тодордог.

\subsection{Design Science Research}
Design Science Research нь мэдээллийн системийн судалгаанд шинэ артефакт бүтээж, түүний хэрэгцээ, зохиомж, үнэлгээг нэгтгэн судлах хандлага юм \cite{hevner_dsr_2004}. Энэхүү дипломын ажил нь зөвхөн онолын судалгаа биш, бодит веб платформ болон RAG dataset pipeline бүтээж байгаа тул DSR-ийн \textit{problem relevance}, \textit{artifact design}, \textit{evaluation}, \textit{research contribution} гэсэн шалгууртай нийцнэ.

\begin{longtable}{|p{0.16\textwidth}|p{0.23\textwidth}|p{0.23\textwidth}|p{0.26\textwidth}|}
  \caption{Аргачлалуудын харьцуулсан шинжилгээ}
  \label{tab:methodology-comparison}\\
  \hline
  \textbf{Аргачлал} & \textbf{Давуу тал} & \textbf{Сул тал} & \textbf{Энэ төсөлд тохирох байдал} \\ \hline
  Waterfall & Баримтжуулалт сайн, үе шат тодорхой & Шаардлага өөрчлөгдөхөд уян хатан бус & Тайлан, шаардлага, зохиомжийг цэгцлэхэд ашиглана. \\ \hline
  Prototype & UX болон шаардлагыг эрт шалгана & Prototype кодыг шууд бүтээгдэхүүн болговол өр үүснэ & Landing, course, lesson, chat UI-д тохиромжтой. \\ \hline
  Spiral & Эрсдэлд суурилсан, AI туршилтад тохиромжтой & Процесс илүү төвөгтэй, туршлага шаарддаг & Dataset, provider, RAG чанарын эрсдэлийг бууруулахад ашиглана. \\ \hline
  Agile/Scrum & Давталттай хөгжүүлэлт, feature backlog, review боломжтой & Бага багт Scrum ceremony бүрэн хэрэгжүүлэх шаардлагагүй & Гол хөгжүүлэлтийн үндсэн аргачлал болгон сонгоно. \\ \hline
  Kanban & Ажлын урсгал ил тод, WIP хянахад сайн & Хугацааны хайрцаг сул байж болно & Дипломын олон төрлийн ажлыг зэрэг удирдахад нэмэлтээр ашиглана. \\ \hline
  Lean UX & Хэрэглэгчийн урсгал хурдан туршихад сайн & Формаль баримтжуулалт бага & Сургалтын интерфейс, chat interaction-д ашиглана. \\ \hline
  DSR & Судалгаа ба бүтээгдэхүүнийг нэгтгэнэ & Үнэлгээний шалгуур тодорхой байх шаардлагатай & Дипломын судалгааны ерөнхий хүрээ болгон ашиглана. \\ \hline
\end{longtable}

\section{Сонгосон арга зүйн загвар}
Дээрх харьцуулалтаас үзэхэд энэ дипломын төсөлд нэг аргачлал дангаар хангалттай биш. Иймээс \textbf{Agile-д суурилсан, Prototype ба Spiral эрсдэлийн хяналтаар баяжуулсан Design Science Research} загварыг сонголоо. Энэ загварын гол санаа нь ажиллаж буй системийг үе шаттай хөгжүүлэхийн зэрэгцээ судалгааны асуудал, артефакт, үнэлгээний логикийг академик байдлаар хадгалах явдал юм.

\begin{table}[htbp]
  \centering
  \caption{Сонгосон арга зүйн бүрэлдэхүүн}
  \label{tab:selected-methodology}
  \begin{tabular}{|p{0.24\textwidth}|p{0.34\textwidth}|p{0.30\textwidth}|}
    \hline
    \textbf{Арга зүйн бүрэлдэхүүн} & \textbf{Төслийн хэрэглээ} & \textbf{Хүлээгдэж буй үр дүн} \\ \hline
    Design Science Research & Судалгааны асуудал тодорхойлох, веб платформ ба RAG туслах артефакт бүтээх & Судалгаа ба бүтээгдэхүүний уялдаа баталгаажна. \\ \hline
    Agile iteration & Feature backlog үүсгэж, модулийг үе шаттай хэрэгжүүлэх & Landing, course, auth, progress, chat зэрэг модулиуд дараалалтай боловсорно. \\ \hline
    Prototype & UX урсгалыг хурдан турших & Course discovery, lesson player, chat drawer ойлгомжтой болно. \\ \hline
    Spiral risk review & AI, dataset, provider, security эрсдэлийг давталт бүрд шалгах & RAG хариултын чанар, fallback, citation-safe бодлого сайжирна. \\ \hline
    Kanban tracking & Thesis, frontend, database, dataset ажлыг зэрэг хянах & Хөгжүүлэлтийн явц ил тод, удирдах боломжтой болно. \\ \hline
  \end{tabular}
\end{table}

\section{Арга зүйн үе шат ба ажлын төлөвлөгөө}
Сонгосон арга зүйг дипломын ажлын бодит үйл ажиллагаанд дараах үе шаттайгаар хэрэгжүүлнэ.

\begin{enumerate}
  \item \textbf{Судалгааны асуудал тодорхойлох:} FL Studio суралцагчдын контентын тархай байдал, ахицын хяналт, AI туслах хэрэгцээг тодорхойлох.
  \item \textbf{Ижил төстэй систем судлах:} онлайн сургалтын платформ, хөгжмийн сургалтын эх сурвалж, AI music систем, RAG chatbot загваруудыг харьцуулах.
  \item \textbf{Шаардлага тодорхойлох:} функциональ ба функциональ бус шаардлагыг backlog хэлбэрээр жагсаах.
  \item \textbf{Prototype боловсруулах:} landing page, course catalog, lesson detail, chat drawer зэрэг хэрэглэгчийн урсгалыг турших.
  \item \textbf{Системийн хэрэгжилт:} frontend, backend/API, database, dataset pipeline, AI provider холболтыг үе шаттай хөгжүүлэх.
  \item \textbf{Эрсдэлийн үнэлгээ:} RAG dataset чанар, source safety, AI хариултын давталт, provider fallback, RLS бодлогыг шалгах.
  \item \textbf{Баримтжуулалт ба дүгнэлт:} судалгааны үр дүн, зохиомж, архитектур, хэрэгжилт, үнэлгээг дипломын тайланд нэгтгэх.
\end{enumerate}

\section{Үнэлгээний арга}
Энэхүү дипломын ажлын үнэлгээ нь зөвхөн код ажиллаж байгаа эсэхийг шалгахаар хязгаарлагдахгүй. Сургалтын платформын хувьд хэрэглэгчийн урсгал, өгөгдлийн бүрэн байдал, AI туслахын хариултын чанар, security policy, өргөтгөх боломж зэрэг олон талаас үнэлэх шаардлагатай.

\begin{table}[htbp]
  \centering
  \caption{Үнэлгээний шалгуур}
  \label{tab:evaluation-criteria}
  \begin{tabular}{|p{0.20\textwidth}|p{0.36\textwidth}|p{0.32\textwidth}|}
    \hline
    \textbf{Шалгуур} & \textbf{Үнэлэх зүйл} & \textbf{Баримтлах арга} \\ \hline
    Функциональ бүрэн байдал & Бүртгэл, курс, хичээл, төлбөр, progress, chatbot ажиллаж байгаа эсэх & Manual scenario test, route бүрийн шалгалт \\ \hline
    Өгөгдлийн загвар & Хүснэгтүүдийн relation, unique constraint, RLS бодлого & SQL schema review, Supabase policy review \\ \hline
    AI туслахын чанар & Retrieval context, prompt, provider response, fallback & Sample question test, source metadata inspection \\ \hline
    UX ойлгомжтой байдал & Course хайх, lesson сонгох, quiz өгөх, chat нээх урсгал & Prototype review, heuristic evaluation \\ \hline
    Баримтжуулалт & Судалгаа, арга зүй, зохиомж, архитектурын тайлбар хангалттай эсэх & Thesis chapter review, diagram/table consistency \\ \hline
  \end{tabular}
\end{table}

\section{Бүлгийн дүгнэлт}
Нэгдүгээр бүлгээр төслийн үүсэл, хөгжлийн үе шат, хэрэглэгчийн хэрэгцээ, функциональ болон функциональ бус шаардлагын эхний түвшний шинжилгээ, ижил төстэй системүүдийн судалгаа, систем хөгжүүлэх аргачлалуудын харьцуулалтыг авч үзлээ. Судалгаанаас харахад энэхүү дипломын төсөл нь зөвхөн AI функцтэй туршилтын сайт бус, харин сургалтын платформ, өгөгдлийн сан, төлбөртэй хандалт, RAG туслахыг нэгтгэсэн иж бүрэн веб систем юм.

Арга зүйн хувьд Waterfall дангаараа уян хатан бус, харин Agile, Prototype, Spiral, Kanban, Design Science Research-ийн хосолсон загвар илүү тохиромжтой гэж дүгнэв. Энэ сонголт нь дараагийн бүлгүүдэд авч үзэх технологийн судалгаа, AI онол, платформын хэрэгжилт, системийн архитектурын логик суурь болно.
